%% Dans ce fichier, les �ventuels accents sont cod�s
%% en ISO-8859-15 (aka latin9)
%% La ligne suivante donne trois exemples 
%% e accent aigu : �, e accent grave : �, e-dans-l'o : �
%%%%%ENTREES%%%%
\begin{tabarg}
\multicolumn{3}{l}{\titretab {Entr\'ees}}\\
\hline
$X$ & Vecteur, \texttt{float} & Param�tre d'entr�e\\
$S_0$ & Coeff., \texttt{float} & Seuil de d�tection d'un pic\\
$S_{mnq}$ & Coeff., \texttt{float} & Seuil de d�tection d'un saut avant ou apr�s une valeur manquante \\
$N_r$ & Coeff., \texttt{float} & Nombre maximum de points invalid�s � remplacer par une interpolation lin�aire\\
\end{tabarg}

\medskip
%%%%%SORTIES%%%%
\begin{tabarg}
\multicolumn{3}{l}{\titretab {Sorties}}\\
\hline
$Flag$ & Vecteur, \texttt{float} & Flag de diagnostic de la correction des pics \\
$X_c$ & Vecteur, \texttt{float} & Param�tre corrig� (m�me unit� que $X$)\\
\end{tabarg}
